\documentclass[12pt]{article}
\usepackage[a4paper,margin=1in]{geometry}
\usepackage{amsmath, amssymb, booktabs, setspace}
\usepackage{enumitem}
\usepackage{graphicx}
\usepackage{siunitx}

\title{Solutions to Practice Problem Set 4\\ \vspace{10pt} \large ECON 101 -- Macroeconomics}
\author{}
\date{Winter 2026}

\begin{document}
\maketitle

%%%%%%%%%%%%%%%%%%%%%%%%%%%%%%%%%%%%%%%%%%%%%%
\section*{Part I: Chapter 6 — Economic Growth and the Financial System}

\begin{enumerate}[label=\textbf{\arabic*.}]

\item \textbf{Difference between economic growth and business cycles}

\textbf{Economic growth} refers to the long-run increase in an economy's ability to produce goods and services. It is usually measured by the growth of \textbf{real GDP} or \textbf{real GDP per person} over time. Economic growth reflects an expansion in the economy's productive capacity.

\textbf{Business cycles} refer to short-run fluctuations in economic activity around the long-run growth trend. These fluctuations include periods of:
\begin{itemize}
    \item expansion,
    \item peak,
    \item recession,
    \item trough.
\end{itemize}

\textbf{Key difference:}
\begin{itemize}
    \item Economic growth is a \textbf{long-run trend}.
    \item Business cycles are \textbf{short-run fluctuations} around that trend.
\end{itemize}

\textbf{Example of economic growth:} Over several decades, Canada's real GDP per person has generally increased because of better technology, more education, and more capital.

\textbf{Example of business cycles:} During the COVID-19 recession, output fell sharply for a short period, but this was a cyclical downturn rather than a permanent decline in long-run growth.

\vspace{0.5em}

\item \textbf{Why labour productivity is crucial for long-run economic growth}

Labour productivity measures output per worker or output per hour worked. It is crucial because living standards rise in the long run only if workers can produce more output.

\textbf{Step 1: Why productivity matters}

If productivity increases, each worker produces more goods and services. This means:
\begin{itemize}
    \item firms can pay higher wages,
    \item more output is available,
    \item living standards improve.
\end{itemize}

\textbf{Step 2: Role of capital deepening}

\textbf{Capital deepening} means an increase in the amount of physical capital per worker, such as more machinery, better equipment, or improved infrastructure.

If workers have better tools and machines, they can produce more in the same amount of time. Therefore, capital deepening raises labour productivity.

\textbf{Step 3: Role of technological progress}

Technological progress means discovering better ways to produce goods and services.

Examples:
\begin{itemize}
    \item improved software,
    \item faster communication systems,
    \item more efficient production methods.
\end{itemize}

Technology allows workers to use existing resources more efficiently, so output rises without requiring the same proportional increase in labour or capital.

\textbf{Conclusion:} Long-run economic growth depends heavily on productivity growth, and productivity growth comes from capital deepening and technological progress.

\vspace{0.5em}

\item \textbf{How financial markets and financial intermediaries contribute to economic growth}

Financial markets and financial intermediaries help move funds from savers to borrowers.

\textbf{Examples:}
\begin{itemize}
    \item \textbf{Financial markets:} stock markets and bond markets.
    \item \textbf{Financial intermediaries:} banks, credit unions, insurance companies.
\end{itemize}

\textbf{Step 1: Channel savings into investment}

Households save money, and firms need funds for investment. Financial institutions connect the two. This allows firms to buy:
\begin{itemize}
    \item machines,
    \item factories,
    \item technology,
    \item research and development.
\end{itemize}

\textbf{Step 2: Improve efficiency}

Financial intermediaries screen borrowers, reduce risk, and lower transaction costs. This makes it easier for resources to go to productive uses.

\textbf{Step 3: Promote growth}

When more savings are transformed into productive investment, the capital stock increases, productivity rises, and long-run growth improves.

\textbf{If financial markets were inefficient:}
\begin{itemize}
    \item savings would not flow easily to firms,
    \item productive investment would fall,
    \item fewer businesses would expand,
    \item capital accumulation and growth would slow.
\end{itemize}

\end{enumerate}

\begin{enumerate}[label=\textbf{\arabic*.}, start=4]

\item \textbf{Loanable Funds Model}

Given:
\begin{itemize}
    \item National saving = \$800 billion,
    \item Investment demand = \$800 billion.
\end{itemize}

Initially, the market is in equilibrium.

\begin{enumerate}
    \item[(a)] \textbf{What happens to national saving?}

If the government runs a deficit of \$200 billion, public saving falls.

National saving is:
\[
S = \text{private saving} + \text{public saving}
\]

A budget deficit means public saving becomes more negative or less positive.

Therefore, national saving falls by \$200 billion.

\[
\boxed{\text{New national saving} = 800 - 200 = 600 \text{ billion}}
\]

\item[(b)] \textbf{How does the supply of loanable funds change?}

The supply of loanable funds equals national saving.

So when national saving falls from 800 to 600, the \textbf{supply curve of loanable funds shifts leftward}.

\textbf{Interpretation:} There are fewer funds available for borrowers in financial markets.

\item[(c)] \textbf{What happens to the equilibrium interest rate?}

When the supply of loanable funds decreases and demand is unchanged:
\begin{itemize}
    \item the equilibrium interest rate rises,
    \item the equilibrium quantity of loanable funds falls.
\end{itemize}

So:
\[
\boxed{\text{The equilibrium interest rate increases}}
\]

\item[(d)] \textbf{Explain crowding out}

\textbf{Crowding out} occurs when government borrowing reduces private investment.

Step-by-step:
\begin{enumerate}
    \item Government deficit increases its need to borrow.
    \item National saving falls.
    \item Supply of loanable funds shifts left.
    \item Interest rates rise.
    \item Higher interest rates make borrowing more expensive for firms.
    \item Firms reduce private investment spending.
\end{enumerate}

Thus:
\[
\boxed{\text{Government borrowing crowds out private investment}}
\]
\end{enumerate}



\item \textbf{Growth Accounting}

Given:
\begin{itemize}
    \item Real GDP growth = 4\%,
    \item Labour input growth = 1\%.
\end{itemize}

\begin{enumerate}
    \item[(a)] \textbf{Calculate labour productivity growth}

Labour productivity growth is approximately:
\[
\text{Labour productivity growth} = \text{Real GDP growth} - \text{Labour input growth}
\]

Substitute the values:
\[
= 4\% - 1\% = 3\%
\]

\[
\boxed{\text{Labour productivity growth} = 3\%}
\]

\item[(b)] \textbf{What factors could explain this growth?}

A 3\% increase in labour productivity can be explained by:
\begin{itemize}
    \item more capital per worker,
    \item technological progress,
    \item improvements in education and training,
    \item better management and organization,
    \item institutional improvements that increase efficiency.
\end{itemize}
\end{enumerate}

\end{enumerate}

\begin{enumerate}[label=\textbf{\arabic*.}, start=6]

\item \textbf{Why some developing countries fail to converge}

Some developing countries do not catch up to rich countries because they face structural barriers.

Three important reasons are:

\textbf{1. Weak institutions}\\
Poor legal systems, corruption, lack of property rights, and political instability discourage investment and innovation.

\textbf{2. Low human capital}\\
If education and health systems are weak, workers are less productive, and growth remains slow.

\textbf{3. Low savings and investment}\\
Without enough saving and investment, countries cannot build the capital stock needed for growth.

Additional reasons may include:
\begin{itemize}
    \item poor infrastructure,
    \item limited access to finance,
    \item trade barriers,
    \item dependence on a narrow range of commodities.
\end{itemize}

\[
\boxed{\text{Failure to converge often reflects deep structural and institutional problems}}
\]

\end{enumerate}

\hrule
\vspace{1em}

%%%%%%%%%%%%%%%%%%%%%%%%%%%%%%%%%%%%%%%%%%%%%%
\section*{Part II: Chapter 8 — Aggregate Expenditure and Output}

\begin{enumerate}[label=\textbf{\arabic*.}, start=7]

\item \textbf{Define MPC and MPS}

\textbf{Marginal Propensity to Consume (MPC)} is the fraction of an additional dollar of income that households spend on consumption.

\[
MPC = \frac{\Delta C}{\Delta Y}
\]

\textbf{Marginal Propensity to Save (MPS)} is the fraction of an additional dollar of income that households save.

\[
MPS = \frac{\Delta S}{\Delta Y}
\]

\textbf{Relationship:}

Since any extra income must either be consumed or saved:
\[
MPC + MPS = 1
\]

So:
\[
\boxed{MPC = 1 - MPS \quad \text{and} \quad MPS = 1 - MPC}
\]

\vspace{0.5em}

\item \textbf{Why the multiplier amplifies changes in spending}

The multiplier effect occurs because one person's spending becomes another person's income.

\textbf{Step-by-step intuition:}
\begin{enumerate}
    \item Government, firms, or households increase spending.
    \item This raises income for someone else.
    \item That person spends part of the extra income.
    \item Their spending becomes income for others.
    \item The process repeats many times.
\end{enumerate}

Because each round creates new spending, the total increase in GDP is larger than the initial change in spending.

\end{enumerate}

\begin{enumerate}[label=\textbf{\arabic*.}, start=9]

\item \textbf{Multiplier Calculation}

Given:
\begin{itemize}
    \item \(MPC = 0.8\),
    \item Increase in government spending \(= \$50\) billion.
\end{itemize}

\begin{enumerate}
    \item[(a)] \textbf{Calculate the multiplier}

The simple spending multiplier is:
\[
k = \frac{1}{1-MPC}
\]

Substitute:
\[
k = \frac{1}{1-0.8} = \frac{1}{0.2} = 5
\]

\[
\boxed{k=5}
\]

\item[(b)] \textbf{Calculate the change in equilibrium GDP}

\[
\Delta Y = k \times \Delta G
\]

\[
\Delta Y = 5 \times 50 = 250
\]

\[
\boxed{\Delta Y = \$250 \text{ billion}}
\]
\end{enumerate}

\item \textbf{Aggregate Expenditure Model}

Given:
\[
C = 200 + 0.75Y, \quad I = 150, \quad G = 250
\]

\begin{enumerate}
    \item[(a)] \textbf{Write the AE equation}

Aggregate expenditure is:
\[
AE = C + I + G
\]

Substitute:
\[
AE = (200 + 0.75Y) + 150 + 250
\]

\[
AE = 600 + 0.75Y
\]

\[
\boxed{AE = 600 + 0.75Y}
\]

\item[(b)] \textbf{Find equilibrium GDP}

Equilibrium occurs where:
\[
Y = AE
\]

So:
\[
Y = 600 + 0.75Y
\]

Subtract \(0.75Y\) from both sides:
\[
Y - 0.75Y = 600
\]

\[
0.25Y = 600
\]

\[
Y = \frac{600}{0.25} = 2400
\]

\[
\boxed{Y = 2400}
\]

\item[(c)] \textbf{Calculate the multiplier}

\[
k = \frac{1}{1-MPC} = \frac{1}{1-0.75} = \frac{1}{0.25} = 4
\]

\[
\boxed{k = 4}
\]

\item[(d)] \textbf{If \(G\) increases by 40, find the new equilibrium GDP}

The change in GDP is:
\[
\Delta Y = k \times \Delta G = 4 \times 40 = 160
\]

New equilibrium GDP:
\[
Y_{\text{new}} = 2400 + 160 = 2560
\]

\[
\boxed{Y_{\text{new}} = 2560}
\]
\end{enumerate}

\end{enumerate}

\begin{enumerate}[label=\textbf{\arabic*.}, start=11]

\newpage

\item \textbf{Recessionary Gap}

Given:
\begin{itemize}
    \item Potential GDP = 3000,
    \item Actual GDP = 2600,
    \item \(MPC = 0.8\).
\end{itemize}

\begin{enumerate}
    \item[(a)] \textbf{Calculate the GDP gap}

\[
\text{GDP gap} = \text{Potential GDP} - \text{Actual GDP}
\]

\[
= 3000 - 2600 = 400
\]

\[
\boxed{\text{Recessionary gap} = 400}
\]

\item[(b)] \textbf{Required increase in government spending}

First calculate the multiplier:
\[
k = \frac{1}{1-MPC} = \frac{1}{1-0.8} = 5
\]

We need:
\[
\Delta Y = 400
\]

Since:
\[
\Delta Y = k \times \Delta G
\]

Then:
\[
400 = 5 \times \Delta G
\]

\[
\Delta G = \frac{400}{5} = 80
\]

\[
\boxed{\Delta G = 80}
\]

\item[(c)] \textbf{Explain the policy intuition}

The economy is producing below potential, so it has a recessionary gap.

An increase in government spending raises aggregate expenditure directly. Through the multiplier process, this initial increase generates additional rounds of income and spending, so equilibrium GDP rises by more than the original increase in \(G\).

Because the multiplier is 5, a \$80 billion increase in government spending closes a \$400 billion gap.

\end{enumerate}

\end{enumerate}

\hrule
\vspace{1em}

%%%%%%%%%%%%%%%%%%%%%%%%%%%%%%%%%%%%%%%%%%%%%%
\section*{Part III: Chapter 9 — Aggregate Demand and Aggregate Supply}

\begin{enumerate}[label=\textbf{\arabic*.}, start=12]

\item \textbf{Why the AD curve slopes downward}

The aggregate demand curve slopes downward for three reasons:

\textbf{1. Wealth effect}\\
When the price level falls, the real value of household wealth rises. People feel wealthier and spend more, so quantity of real GDP demanded increases.

\textbf{2. Interest rate effect}\\
When the price level falls, households and firms need less money for transactions. Interest rates fall, which encourages consumption and investment spending.

\textbf{3. International trade effect}\\
When the domestic price level falls relative to foreign prices, domestic goods become relatively cheaper. Exports rise and imports fall, so net exports increase.

Thus:
\[
\boxed{\text{A lower price level increases the quantity of real GDP demanded}}
\]

\vspace{0.5em}

\item \textbf{Distinguish between SRAS and LRAS}

\textbf{Short-run aggregate supply (SRAS):}
\begin{itemize}
    \item slopes upward,
    \item in the short run, some input prices are sticky,
    \item firms respond to higher output prices by increasing production.
\end{itemize}

\textbf{Long-run aggregate supply (LRAS):}
\begin{itemize}
    \item is vertical,
    \item in the long run, output depends on labour, capital, technology, and institutions,
    \item the price level does not affect potential GDP.
\end{itemize}

So:
\[
\boxed{\text{SRAS reflects short-run price-output responses, while LRAS reflects potential output}}
\]

\vspace{0.5em}

\item \textbf{What causes shifts in AD, SRAS, and LRAS?}

\textbf{AD shifts because of changes in:}
\begin{itemize}
    \item consumption,
    \item investment,
    \item government spending,
    \item net exports.
\end{itemize}

\textbf{Example:} An increase in government spending shifts AD right.

\textbf{SRAS shifts because of changes in:}
\begin{itemize}
    \item input prices,
    \item productivity,
    \item supply shocks.
\end{itemize}

\textbf{Example:} A rise in oil prices shifts SRAS left.

\textbf{LRAS shifts because of changes in:}
\begin{itemize}
    \item labour force,
    \item capital stock,
    \item technology,
    \item institutions.
\end{itemize}

\textbf{Example:} Technological improvement shifts LRAS right.

\end{enumerate}

\begin{enumerate}[label=\textbf{\arabic*.}, start=15]

\item \textbf{Demand Shock}

Consumer confidence falls.

\begin{enumerate}
    \item[(a)] \textbf{What happens to AD?}

Lower consumer confidence reduces household spending, especially on durable goods. Therefore:
\[
\boxed{\text{AD shifts left}}
\]

\item[(b)] \textbf{Short-run effects}

With AD shifting left:
\begin{itemize}
    \item \textbf{output falls},
    \item \textbf{price level falls} (or inflation slows),
    \item \textbf{unemployment rises}.
\end{itemize}

Reason: firms sell less output, reduce production, and hire fewer workers.
\end{enumerate}

\newpage

\item \textbf{Supply Shock}

Oil prices increase significantly.

\begin{enumerate}
    \item[(a)] \textbf{What happens to SRAS?}

Oil is an important production input. Higher oil prices raise production costs.

Therefore:
\[
\boxed{\text{SRAS shifts left}}
\]

\item[(b)] \textbf{What happens to inflation and output?}

A leftward shift of SRAS causes:
\begin{itemize}
    \item \textbf{higher price level / higher inflation},
    \item \textbf{lower output}.
\end{itemize}

\item[(c)] \textbf{What type of inflation is this?}

This is:
\[
\boxed{\text{cost-push inflation}}
\]

because rising production costs push the price level up.
\end{enumerate}

\end{enumerate}

\begin{enumerate}[label=\textbf{\arabic*.}, start=17]

\item Given:
\begin{itemize}
    \item Potential GDP = 2500,
    \item Actual GDP = 2700,
    \item Inflation is rising.
\end{itemize}

\begin{enumerate}
    \item[(a)] \textbf{Inflationary or recessionary gap?}

Since actual GDP exceeds potential GDP:
\[
2700 > 2500
\]

\[
\text{Gap} = 2700 - 2500 = 200
\]

So this is an:
\[
\boxed{\text{inflationary gap of 200}}
\]

\item[(b)] \textbf{What policy should be used?}

Because the economy is overheating and inflation is rising, policymakers should use:
\begin{itemize}
    \item contractionary fiscal policy, or
    \item contractionary monetary policy.
\end{itemize}

\[
\boxed{\text{A contractionary policy should be used}}
\]

\item[(c)] \textbf{Adjustment process back to equilibrium}

Step-by-step:
\begin{enumerate}
    \item The economy is above potential GDP.
    \item Labour markets become tight and wages rise.
    \item Rising wages increase firms' costs.
    \item SRAS shifts left over time.
    \item Output falls back toward potential GDP.
    \item Price level rises during the adjustment.
\end{enumerate}

If policymakers act directly, AD may shift left sooner, reducing inflationary pressure.
\end{enumerate}

\end{enumerate}

\begin{enumerate}[label=\textbf{\arabic*.}, start=18]

\item The economy experiences:
\begin{itemize}
    \item decline in investment,
    \item rising interest rates,
    \item falling GDP.
\end{itemize}

\begin{enumerate}
    \item[(a)] \textbf{Explain using AE and AD--AS}

\textbf{AE model:}
\begin{itemize}
    \item Investment is a component of aggregate expenditure.
    \item A decline in investment shifts the AE line downward.
    \item Through the multiplier process, equilibrium GDP falls by more than the initial decline.
\end{itemize}

\textbf{AD--AS model:}
\begin{itemize}
    \item A decline in investment reduces aggregate demand.
    \item AD shifts left.
    \item In the short run, output falls and the price level may also fall or rise more slowly.
\end{itemize}

\item[(b)] \textbf{Role of the financial system}

The financial system affects the downturn because:
\begin{itemize}
    \item rising interest rates make borrowing more expensive,
    \item firms reduce investment spending,
    \item households may cut back on major purchases,
    \item banks may lend less if financial conditions worsen.
\end{itemize}

Thus, weakness in the financial system can amplify the recession.

\item[(c)] \textbf{One fiscal and one monetary policy}

\textbf{Fiscal policy:} Increase government spending or reduce taxes.\\
This raises aggregate expenditure and shifts AD right.

\textbf{Monetary policy:} Lower the policy interest rate.\\
This reduces borrowing costs, encourages investment and consumption, and shifts AD right.

Both policies help increase GDP and reduce unemployment.
\end{enumerate}

\end{enumerate}

\begin{enumerate}[label=\textbf{\arabic*.}, start=19]

\item Government spending increases by \$100, \(MPC=0.75\), but interest rates rise because of deficit financing.

First compute the simple multiplier:
\[
k = \frac{1}{1-MPC} = \frac{1}{1-0.75} = 4
\]

Predicted increase in GDP without crowding out:
\[
\Delta Y = 4 \times 100 = 400
\]

So the simple multiplier predicts:
\[
\boxed{\Delta Y = 400}
\]

\textbf{Why actual GDP may rise by less than 400:}

If the government finances the increase in spending by borrowing, this may reduce national saving and raise interest rates.

Higher interest rates can:
\begin{itemize}
    \item reduce private investment,
    \item reduce interest-sensitive consumption.
\end{itemize}

So while government spending rises, some private spending falls.

This is \textbf{crowding out}, and it weakens the multiplier effect.

Therefore:
\[
\boxed{\text{Actual GDP rises by less than the simple multiplier predicts}}
\]

\end{enumerate}

\hrule
\vspace{1em}

\begin{center}
\textit{End of Solution Key}
\end{center}

\end{document}